Large language models generate fluent but factually incorrect text, yet detecting these hallucinations remains challenging because existing uncertainty methods have been developed and evaluated in isolation. We present the first systematic comparison of entropy-based and consistency-based uncertainty signals for hallucination detection on the same factuality benchmark. Our key finding is that these signals are orthogonal: token entropy and N-sample consistency share only 5\% of variance ($r = 0.228$), capturing fundamentally different failure modes---entropy reflects epistemic uncertainty while consistency reflects generation stability. On 18\% of questions where the methods disagree, each achieves strong predictive performance (AUROC $> 0.75$) on its winning subset, demonstrating genuine complementary detection capability. These results establish the theoretical foundation for multi-signal hallucination detection and motivate hybrid approaches that combine orthogonal uncertainty signals for more robust error detection in deployed LLM systems.
